\documentclass[10pt]{article}
\usepackage[utf8]{inputenc}
\usepackage{amsmath}
\usepackage{geometry}
\usepackage{enumitem}
\geometry{a4paper, margin=0.5in}

\setlist[enumerate]{itemsep=1pt, topsep=2pt, parsep=0pt}

\begin{document}

\begin{center}
    \textbf{\large EXAMEN PARCIAL DE FÍSICA BÁSICA (FIS-0001) -- TEMA A} \\
    \textbf{INSTITUTO TECNOLÓGICO SUPERIOR COMUNITARIO (ITSC)} \\
    \textbf{Docente: Lic. Carlos Ramírez \hfill Calificación: \rule{1.5cm}{0.1pt} / 10 Puntos}
\end{center}

\vspace{0.1cm}
\noindent Nombre: \hrulefill \ Matrícula: \hrulefill \ Sección: \hrulefill \ Fecha: \hrulefill

\vspace{0.2cm}
\begin{center}
    \framebox[\textwidth]{\parbox{0.95\textwidth}{\centering \small Fórmulas de apoyo:
        $d = v \cdot t \quad \vec{R} = (A_x + B_x, A_y + B_y) \quad A_y = A \operatorname{sen}\theta \quad A = \sqrt{A_x^2 + A_y^2} \quad t = \frac{d}{v} \quad \tan\theta = \frac{A_y}{A_x} \quad v = \frac{d}{t} \quad A_x = A \cos\theta$
    }}
\end{center}

\vspace{0.2cm}
\noindent \textbf{PARTE I: CUESTIONARIO TEÓRICO (4 Puntos -- 0.5 pto c/u)} \\
\textit{Instrucciones: Encierre en un círculo la letra que represente la opción correcta para cada enunciado.}

\begin{enumerate}[leftmargin=*]
    \item ¿Cuál es el objetivo principal de expresar una cantidad física en notación científica? \\
    a) Eliminar los números decimales para facilitar la lectura. \\
    b) Escribir cantidades muy grandes o muy pequeñas de forma compacta utilizando potencias de base 10. \\
    c) Aumentar la precisión matemática de una medición experimental.

    \item ¿Cuántas cifras significativas tiene la cantidad medida $0.007030\text{ A}$? \\
    a) Tres cifras significativas \hfill b) Cuatro cifras significativas \hfill c) Seis cifras significativas \hfill \mbox{}

    \item Al sumar las mediciones $12.3\text{ g}$ y $1.45\text{ g}$, el resultado final expresado con el número correcto de cifras significativas es: \\
    a) $13.75\text{ g}$ \hfill b) $13.8\text{ g}$ \hfill c) $14\text{ g}$ \hfill \mbox{}

    \item ¿Cuál de las siguientes cantidades físicas es una magnitud vectorial? \\
    a) La temperatura \hfill b) La masa \hfill c) La fuerza \hfill \mbox{}

    \item Un vector se define físicamente como un segmento de recta dirigido que posee: \\
    a) Solamente módulo y sentido. \\
    b) Módulo, dirección y sentido. \\
    c) Módulo y coordenadas cartesianas únicamente.

    \item Si dos vectores en el plano cartesiano son perpendiculares entre sí, su suma vectorial dará como resultado: \\
    a) Un vector resultante cuyo módulo es igual a la suma aritmética de sus módulos. \\
    b) Un vector resultante cuyo módulo se calcula aplicando el Teorema de Pitágoras. \\
    c) Un vector nulo (cero).

    \item La línea recta que une directamente el punto de partida (inicial) con el punto de llegada (final) de un móvil se denomina: \\
    a) Distancia recorrida \hfill b) Desplazamiento resultante \hfill c) Trayectoria real \hfill \mbox{}

    \item ¿Cuál es la principal característica del Movimiento Rectilíneo Uniforme (M.R.U.)? \\
    a) Su velocidad cambia de forma constante y uniforme con el tiempo. \\
    b) Su aceleración es constante y diferente de cero. \\
    c) Su velocidad es constante en módulo, dirección y sentido, por lo que su aceleración es cero.
\end{enumerate}

\vspace{0.3cm}
\noindent \textbf{PARTE II: EJERCICIOS PRÁCTICOS (6 Puntos -- 1.5 ptos c/u)} \\
\textit{Instrucciones: Resuelva de forma ordenada cada ejercicio. Los problemas deben ser justificados en la parte de atrás de esta hoja. Escriba el resultado final en la línea correspondiente.}

\begin{enumerate}[leftmargin=*]
    \item Realice la siguiente operación y exprese el resultado final en notación científica: \\
    $(4.00 \times 10^4\text{ m/s}) \times (2.00 \times 10^{-2}\text{ s})$ \\
    \textbf{Respuesta:} Resultado = \rule{5cm}{0.1pt} m

    \item Efectúe la siguiente suma de mediciones y exprese el resultado final redondeado al número correcto de cifras significativas: \\
    $12.35\text{ kg} + 4.2\text{ kg} + 0.156\text{ kg}$ \\
    \textbf{Respuesta:} Suma = \rule{5cm}{0.1pt} kg

    \item Un vector fuerza $\vec{F}$ tiene una magnitud de $20.00\text{ N}$ y forma un ángulo de $30.00^\circ$ con el eje horizontal positivo. Determine los valores de sus componentes rectangulares $F_x$ y $F_y$. (Datos: $\cos 30.00^\circ \approx 0.87$, $\operatorname{sen} 30.00^\circ = 0.50$) \\
    \textbf{Respuesta:} $F_x = $ \rule{3cm}{0.1pt} N \quad | \quad $F_y = $ \rule{3cm}{0.1pt} N

    \item Un automóvil viaja por una autopista recta manteniendo una velocidad constante de $25.00\text{ m/s}$. Calcule la distancia en metros que logrará recorrer al cabo de un intervalo de tiempo de $2.00\text{ minutos}$ ($120.00\text{ segundos}$). \\
    \textbf{Respuesta:} Distancia = \rule{5cm}{0.1pt} m
\end{enumerate}

% CLAVE DE RESPUESTAS PARA CONTROL DEL DOCENTE (SÓLO CÓDIGO FUENTE):
% PARTE I (Teoría):
% 1. b
% 2. b (Cifras significativas: 7, 0, 3, 0)
% 3. b (Se redondea a un decimal debido a 12.3: 13.75 -> 13.8)
% 4. c (La fuerza es vectorial; temperatura y masa son escalares)
% 5. b
% 6. b
% 7. b
% 8. c
% 
% PARTE II (Ejercicios):
% 1. 8.00 \times 10^2 m (o 800.00 m)
% 2. 16.7 kg (Suma matemática da 16.706, pero se limita a un decimal por el sumando "4.2")
% 3. F_x = 17.40 N (20.00 * 0.87) | F_y = 10.00 N (20.00 * 0.50)
% 4. 3000.00 m o 3.00 \times 10^3 m (d = v * t = 25.00 * 120.00)

\end{document}
